\section{弱 \texorpdfstring{$J$}{J}-全纯映射的边界正则性}
\stfhIndex{I065}{bian^jie^zheng^ze^xing}{边界正则性（boundary regularity）}
\label{sec:boundary-regularity-weak-j-holomorphic}
\stfhIndex{I304}{ruo^j^quan^chun}{弱 $J$-全纯（weakly $J$-holomorphic）}
\stfhIndex{I302}{ruo^nuo^yi^man^bian^jie^tiao^jian}{弱诺伊曼边界条件（weak Neumann boundary condition）}

闭曲面 $\Sigma$ 上弱 $J$-全纯映射的讨论较易处理，并且可纳入带边曲面的研究。因此，以下讨论集中于有边界情形。

设 $(N,J,g)$ 为殆埃尔米特流形，$R\subset N$ 是关于 $J$ 的 $n$ 维全实子流形。我们引入满足全实边界条件
\begin{equation}
 \bar\partial_Ju=0,\qquad u(\partial\Sigma)\subset R.
 \label{eq:ch8-weak-j-holomorphic-system}
\end{equation}
的 $J$-全纯映射之\emph{弱解}概念。为此，需要使用引理 8.3.1 所选的埃尔米特度量
$g_0$，使 $TR\perp_{g_0}JTR$。由纳什嵌入定理，可假设 $(N,g_0)$ 等距嵌入
$\mathbb R^N$，并把映射 $u$ 及沿 $u$ 的向量场 $\xi$ 看作从 $\Sigma$ 到
$\mathbb R^N$ 的映射。特别地，可以把到 $M$ 的 $L^1$ 映射当作到 $\mathbb R^N$
的 $L^1$ 映射，并要求它在定义域上几乎处处取值于 $M$。与前几节不同，以下使用
$g_0$ 的列维--奇维塔联络。明白这些约定后，受 (Ye94) 中定义的启发，作如下定义。

\begin{definition}\label{def:ch8-weak-j-holomorphic-neumann}
称 $W^{1,2}$ 映射 $u:(\Sigma,j)\to(N,J)$ 是 \eqref{eq:ch8-weak-j-holomorphic-system} 关于 $g_0$ 的\emph{弱解}，如果它在 $\partial\Sigma$ 上为 $L^1$ 且几乎处处映入 $R$，并满足：
\begin{enumerate}[label=(\arabic*)]
\item $\bar\partial_Ju=0$ 在 $\Sigma$ 上几乎处处成立；
\item
\begin{equation}
 \int_\Sigma\langle Jdu,\nabla X\rangle=0
 \label{eq:ch8-weak-neumann-condition}
\end{equation}
对每个沿 $u$ 的光滑向量场 $X$ 成立；这里 $z\in\partial\Sigma$ 时
$X(z)\perp_{g_0}R$，且 $\int_{\partial\Sigma}\langle\nabla X,\nabla X\rangle<\infty$。
\end{enumerate}
也称 $u$ 为在 $R$ 上满足诺伊曼边界条件（关于 $g_0$）的\emph{弱 $J$-全纯映射}。
\end{definition}

\begin{lemma}\label{lem:ch8-weak-implies-classical-neumann}
任何满足 \eqref{eq:ch8-weak-neumann-condition} 的 $C^1$ 解都在经典意义下满足诺伊曼边界条件。
\end{lemma}
\begin{proof}
先设 $u\in C^2$。在极坐标中用
$u_r+Jr^{-1}u_\theta=0$。若 $X$ 在 $\mathring D$ 内紧支，并且在边界上关于 $g_0$ 垂直于 $R$，则由斯托克斯公式、分部积分及列维--奇维塔联络无挠，
\begin{align*}
 \int_D\langle Jdu,\nabla X\rangle
 &=\int_D\left(\left\langle J\frac{\partial u}{\partial r},
 \frac{DX}{\partial r}\right\rangle
 +\frac1{r^2}\left\langle J\frac{\partial u}{\partial\theta},
 \frac{DX}{\partial\theta}\right\rangle\right)r\,dr\,d\theta\\
 &=\int_D\left(\left\langle\frac1r\frac{\partial u}{\partial\theta},
 \frac{DX}{\partial r}\right\rangle
 -\left\langle\frac{\partial u}{\partial r},
 \frac1r\frac{DX}{\partial\theta}\right\rangle\right)r\,dr\,d\theta\\
 &=-\int_D\left\langle\frac D{\partial r}\frac{\partial u}{\partial\theta},X\right\rangle\,dr\,d\theta
 +\int_D\left\langle\frac D{\partial\theta}\frac{\partial u}{\partial r},X\right\rangle\,dr\,d\theta
 -\int_{\partial D}\left\langle\frac{\partial u}{\partial\theta},X\right\rangle d\theta\\
 &=-\int_{\partial D}\left\langle\frac{\partial u}{\partial\theta},X\right\rangle d\theta
 =\int_{\partial D}\left\langle\frac{\partial u}{\partial r},JX\right\rangle d\theta,
\end{align*}
其中第三个等号使用分部积分；第四个等号使用列维--奇维塔联络无挠所给出的恒等式
\[
 \frac D{\partial\theta}\frac{\partial u}{\partial r}
 =\frac D{\partial r}\frac{\partial u}{\partial\theta}.
\]
因此得到
\[
 -\int_{\partial D}\left\langle\frac{\partial u}{\partial r},JX\right\rangle d\theta=0.
\]
这对所有光滑 $X\perp TL$ 成立，等价地，$JX$ 沿 $\partial D$ 切于 $L$，故
$\partial u/\partial r\perp L$。对 $C^1$ 映射以 $C^2$ 映射作 $C^1$ 逼近并取极限。这就完成了证明。
\end{proof}
\begin{exercise}\label{ex:ch8-fill-limiting-argument}
补全上述证明末尾的极限论证。
\end{exercise}

为把弱边界条件 \eqref{eq:ch8-weak-neumann-condition} 的含义转化为
$\mathbb R^{2n}$ 中更标准的分布边界条件，先用沿 $R$ 的列维--奇维塔联络活动标架
写出带诺伊曼边界条件的 $J$-全纯方程；类似计算见托马斯·帕克（Thomas H. Parker）与乔恩·沃尔夫森（Jon Wolfson）
(PW93) 的附录。取 $N$ 的正交标架
\begin{equation}
 \{e_1,\ldots,e_n,f_1,\ldots,f_n\},\qquad f_j=Je_j,
 \label{eq:ch8-adapted-real-frame}
\end{equation}
并令 $u_j=2^{-1/2}(e_j-if_j)$，则
\begin{equation}
 \operatorname{Re}u_j|_R=2^{-1/2}e_j.
 \label{eq:ch8-real-part-unitary-frame}
\end{equation}
这些 $u_j$ 构成 $T^{(1,0)}N$ 的酉标架。令
\begin{equation}
 \{\alpha_1,\ldots,\alpha_n,\beta_1,\ldots,\beta_n\}
 \label{eq:ch8-dual-real-frame}
\end{equation}
为 \eqref{eq:ch8-adapted-real-frame} 的对偶标架。复值 $1$-形式
$\theta^j=2^{-1/2}(\alpha^j+i\beta^j)$ 构成
$(T_\mathbb C^*M)^{(1,0)}$ 的酉标架，并与 $\{u_1,\ldots,u_n\}$ 对偶。
殆埃尔米特联络由局部矩阵值复 $1$-形式
$\omega=(\omega^i_j)_{i,j=1,2,\ldots,n}$ 表示。取
$\Sigma$ 上正交余标架 $\{\eta^1,\eta^2\}$，令 $\phi=\eta^1+i\eta^2$；结构方程为
\begin{equation}
 d\phi=-i\rho\wedge\phi=-i\rho_{\bar\phi}\bar\phi\wedge\phi.
 \label{eq:ch8-domain-structure-equation}
\end{equation}
$\rho=\rho_\phi\phi+\rho_{\bar\phi}\bar\phi$ 是余标架的列维--奇维塔联络形式，而高斯曲率由第二结构方程
\[
 d\rho=-\frac i2K\phi\wedge\bar\phi
\]
定义。
$J$-全纯性等价于
\begin{equation}
 u^*\theta^j=a^j\phi,
 \label{eq:ch8-pullback-unitary-coframe}
\end{equation}
且
\[
 |du|^2=2\sum_{j=1}^n|a^j|^2.
\]
由于假设 $du\in L^2$，各 $a^j$ 也在 $L^2$。利用 $u$ 的连续性，并在给定边界点 $u(z)\in Y$ 附近取足够小的邻域，可假设标架及 $a^j$ 在该邻域全局定义。由第 7 章的定义，
\[
 (\partial u)'=\sum_i u^*\theta^i\cdot u_i
 =\sum_i(a^i\phi)\cdot u_i.
\]
由定理 7.3.4 的 $\bar\partial^\nabla(\partial u)'=0$ 得
\begin{equation}
 0=\sum_i\bar\partial(a^i\phi)u_i+\sum_j a^j\phi\wedge\sum_i a^i\omega^i_j u_i,
 \label{eq:ch8-moving-frame-preliminary}
\end{equation}
用 \eqref{eq:ch8-domain-structure-equation} 计算
\[
 \bar\partial(a^i\phi)=\bar\partial a^i\wedge\phi
 -i\rho_{\bar\phi}\bar\phi\wedge\phi.
\]
依标准记号记
\[
 \nabla'a^i:=da^i\cdot\phi,\qquad
 \nabla''a^i:=da^i\cdot\bar\phi,\qquad
 \omega^i_j=\omega^i_{jk}\theta^k+\omega^i_{j\bar k}\bar\theta^k.
\]
从而
\begin{equation}
 \nabla''a^i-i\rho_{\bar\phi}a^i-\omega^i_{j\bar\ell}a^j\bar a^\ell=0.
 \label{eq:ch8-moving-frame-equation}
\end{equation}
另一方面，由 \eqref{eq:ch8-real-part-unitary-frame}，
$\operatorname{Re}\theta^j|_{TR}$ 取实值；由诺伊曼边界条件，
\begin{equation}
 du(\partial/\partial\nu)\in\operatorname{span}\{f_j\},
 \label{eq:ch8-neumann-moving-frame}
\end{equation}
另一方面，由 $du$ 为实 $1$-形式，
\[
 du=\sum_j u^*\theta^j\cdot u_j+\sum_k u^*\bar\theta^k\cdot\bar u_k
 =\sum_j(a^j\phi\cdot u_j+\bar a^j\bar\phi\cdot\bar u_j).
\]
从而
\[
 du\left(\frac\partial{\partial\nu}\right)
 =\sum_j\bigl(a^j\,\phi\bigl(\frac\partial{\partial\nu}\bigr)\cdot u_j
   +\bar a^j\,\bar\phi\bigl(\frac\partial{\partial\nu}\bigr)\cdot\bar u_j\bigr).
\]
取等温坐标 $(r,\varphi)$，使 $\partial/\partial\nu=\partial/\partial r$ 且
$\phi=\mu(dr+i\,d\varphi)$、$\mu>0$，则 $\phi(\partial/\partial\nu)=\bar\phi(\partial/\partial\nu)=\mu$
均为实函数。于是
\[
 du\left(\frac\partial{\partial\nu}\right)
 =\mu\sum_j(a^ju_j+\bar a^j\bar u_j)
 =\mu\sum_j\bigl((a^j+\bar a^j)e_j-i(a^j-\bar a^j)f_j\bigr).
\]
所以条件 \eqref{eq:ch8-neumann-moving-frame} 等价于
$a^j+\bar a^j=0$，即各 $a^j$ 在 $\partial\Sigma$ 上纯虚。

上述计算概括如下。
\begin{proposition}\label{prop:ch8-adapted-frame-boundary-system}
设 $\{\theta^j\}$ 是如上的适合于 $R$ 的活动标架，并设 $u$ 是沿
$\partial\Sigma$ 在 $R$ 上满足诺伊曼边界条件的光滑 $J$-全纯映射。
对每个 $j=1,\ldots,n$，把 $u^*\theta^j=a^j\phi$ 记为复值函数 $a^j$。则
$\bar\partial^\nabla(\partial u)=0$ 等价于
\begin{equation}
 \begin{cases}
 \nabla''a^i-i\rho_{\bar\phi}a^i-\omega^i_{j\bar\ell}a^j\bar a^\ell=0,\\
 a^i\text{ 在边界上为纯虚数},\quad i=1,\ldots,n.
 \end{cases}
 \label{eq:ch8-adapted-frame-boundary-value}
\end{equation}
\end{proposition}

下面这个具有诺伊曼边界条件的弱 $J$-全纯映射正则性定理是本节主定理。
\begin{theorem}\label{thm:ch8-weak-boundary-regularity}
任何在 $\partial\Sigma$ 上沿 $R$ 满足诺伊曼边界条件的弱 $J$-全纯映射
$u:\Sigma\to M$，若属于某个 $p\geq4$ 的 $W^{1,p}$，或属于某个
$\epsilon>0$ 的 $C^\epsilon$，则光滑。
\end{theorem}
\begin{proof}
先注意，$p>2$ 且 $0\leq\epsilon<1-2/p$ 时，索伯列夫嵌入
$W^{1,p}\hookrightarrow C^\epsilon$ 蕴含 $u$ 是带弱边界条件
\eqref{eq:ch8-weak-neumann-condition} 的 \eqref{eq:ch8-weak-j-holomorphic-system}
之 $C^\epsilon$ 解。由于 $u$ 连续，可以局部研究正则性。内部正则性较容易，以下只论边界。

设 $\{\theta^j\}$ 是如上的适合于 $R$ 的 $(T_\mathbb C^*M)^{(1,0)}$
之活动标架。用舍恩--乌伦贝克强逼近定理取光滑
\stfhIndex{I243}{she^en^wu^lun^bei^ke^qiang^bi^jin^ding^li}{舍恩--乌伦贝克强逼近定理（Schoen–Uhlenbeck strong approximation theorem）}
$u_\alpha\to u$ 于 $W^{1,p}$，于是
\[
 \int_\Sigma|\bar\partial_Ju_\alpha|^2\to0,
\qquad
 \int_\Sigma\langle\partial u_\alpha,(\bar\partial^\nabla)^*X\rangle
 \to\int_\Sigma\langle\partial u,(\bar\partial^\nabla)^*X\rangle=0,
\]
对沿 $u$ 且 $X(z)$ 切于 $Y$、并满足
$\int_{\partial\Sigma}\langle\nabla X,\nabla X\rangle<\infty$ 的向量场 $X$ 成立。写
$u_\alpha^*\theta^j=a_\alpha^j\phi+b_\alpha^j\bar\phi$；则
$b_\alpha^j\to0$，$a_\alpha^j\to a^j$ 于 $L^p$，而 $a$ 弱满足
\eqref{eq:ch8-adapted-frame-boundary-value}。其低阶项
\begin{equation}
 -i\rho_{\bar\phi}a^i-\omega^i_{j\bar\ell}a^j\bar a^\ell
 \label{eq:ch8-moving-frame-lower-order}
\end{equation}
属于 $L^{p/2}$；因 $p\geq4$，有 $p/2\geq2$。还要注意，把
$\nabla''=\bar\partial+\Gamma$ 看作系数在 $C^\epsilon$ 中的零阶扰动，方程具有形式
\[
 \begin{cases}
 \bar\partial a^i=f^i,\\
 a^j\text{ 在边界上为纯虚数},
 \end{cases}
 \qquad f^i\in L^p.
\]
这是主符号与 $\bar\partial$ 相同的一阶椭圆边值问题。标准正则性结果
(Mo66)（或定理 8.3.5 的特例）给出
\[
 \|da\|_{1,p,K}\leq C_2\|f\|_{0,p,D}
 \leq C_3\bigl(\|a\|_{0,p,D}+\|a\|_{0,p/2,D}\bigr),
\]
其中结合使用了 \eqref{eq:ch8-adapted-frame-boundary-value}。因此
$du\in W^{2,p}\hookrightarrow C^\epsilon$，从而 $u\in C^{1,\epsilon}$。
把 \eqref{eq:ch8-adapted-frame-boundary-value} 看作系数为赫尔德连续函数的
$a^j$ 线性方程，再作标准 bootstrap，便得到更高正则性。这就完成了证明。
\end{proof}

由此，局部正则性的研究归结为取得 $p\geq4$ 时 $J$-全纯映射的
$W^{1,p}$ 估计，特别是有限面积或有限调和能量映射的这种估计。
